\section{The Harness-Hacker Protocol}

\label{sec:protocol}

\subsection{Protocol Overview}
\label{sec:protocol-overview}

The Harness-Hacker Protocol specifies a six-state machine governing
self-modification attempts. Figure~\ref{fig:state-machine} summarizes the
transitions; Algorithm~\ref{alg:hhp} provides the complete pseudocode.

\begin{figure}[H]
\centering
\setlength{\fboxsep}{8pt}
\fbox{%
\begin{minipage}{0.9\textwidth}
\centering
\textbf{HHP State Machine}

\vspace{6pt}
$\texttt{DETECT} \;\xrightarrow{\text{scope classified}}\; \texttt{ISOLATE}
\;\xrightarrow{\text{worktree created}}\; \texttt{MODIFY}$

\vspace{4pt}
$\texttt{MODIFY} \;\xrightarrow{\text{diff produced}}\; \texttt{TEST}
\;\xrightarrow{\Gamma_k = \text{pass}}\; \texttt{GATE}
\;\xrightarrow{\text{approved}}\; \texttt{DEPLOY}$

\vspace{4pt}
$\texttt{TEST} \;\xrightarrow{\Gamma_k = \text{fail}}\; \texttt{MODIFY}$
\quad (max 3 fix attempts)

\vspace{4pt}
$\texttt{GATE} \;\xrightarrow{\text{rejected}}\; \texttt{DETECT}$
\quad (friction event archived, no modification)
\end{minipage}
}
\caption{The six-state HHP state machine. Self-loops between MODIFY and TEST
are bounded at three attempts before the hack cycle is abandoned.}
\label{fig:state-machine}
\end{figure}

\subsection{State Transition Specifications}
\label{sec:transitions}

\paragraph{DETECT.}
\textbf{Precondition:} A friction event $\mathcal{F} = (\tau, e, n)$ has been
observed. \\
\textbf{Action:} Classify the friction into a modification scope
$k \in \{1, 2, 3, 4\}$ using Algorithm~\ref{alg:classify}. \\
\textbf{Postcondition:} Scope $k$ is determined and the modification session
budget is checked: $\text{SessionHacks} < 2$.

\paragraph{ISOLATE.}
\textbf{Precondition:} Scope $k$ and target component $\mathcal{R}$ are
determined. \\
\textbf{Action:} For $k \in \{2, 3, 4\}$: create a git worktree
\texttt{git worktree add /tmp/hack-\$ID origin/main}. For $k = 1$: allocate
an in-process sandbox. \\
\textbf{Postcondition:} Isolation environment $B$ exists and satisfies
Definition~\ref{def:isolation}. The running system state $\Sigma$ is unchanged.

\paragraph{MODIFY.}
\textbf{Precondition:} Isolation environment $B$ is established. \\
\textbf{Action:} The agent writes modifications within $B$, producing diff
$\Delta$. For Tier 1, this is an executable function. For Tiers 2--4, this is
git-tracked source changes. \\
\textbf{Postcondition:} $\Delta$ is complete and syntactically valid. Fix
attempt counter is incremented if this is a re-entry from TEST.

\paragraph{TEST.}
\textbf{Precondition:} $\Delta$ is complete. \\
\textbf{Action:} Apply validation gate $\Gamma_k(\Delta, \Sigma)$. \\
\textbf{Postcondition:} If $\Gamma_k = \text{pass}$, transition to GATE.
If $\Gamma_k = \text{fail}$ and fix-attempts $< 3$, transition back to MODIFY
with the failure details appended to context. If fix-attempts $= 3$, abandon
the hack cycle, archive the friction event, and return to normal operation.

\paragraph{GATE.}
\textbf{Precondition:} $\Gamma_k(\Delta, \Sigma) = \text{pass}$. \\
\textbf{Action:} Submit $\Delta$ to the tier-appropriate gate:
Tier 1: register immediately (gate is $\Gamma_1$, already passed in TEST).
Tiers 2--3: invoke automated reviewer.
Tier 4: open pull request on target repository. \\
\textbf{Postcondition:} Deployment decision is received ($\text{approved}$ or
$\text{rejected}$).

\paragraph{DEPLOY.}
\textbf{Precondition:} Gate decision is $\text{approved}$. \\
\textbf{Action:} Apply $D_k(\Delta)$ (tier-appropriate deployment). \\
\textbf{Postcondition:} Modification is live. Session hack counter is
incremented. Worktree is cleaned up: \texttt{git worktree remove --force \$B}.

\subsection{The HARNESS-HACK Algorithm}
\label{sec:algorithm}

\begin{algorithm}[H]
\caption{HARNESS-HACK: Safe Agent Self-Modification}
\label{alg:hhp}
\begin{algorithmic}[1]
\Require Friction event $\mathcal{F} = (\tau, e, n)$, session hack counter $H$
\Ensure Modification deployed or friction archived; $H$ updated
\If{$H \geq 2$}
    \State \textbf{Archive} $\mathcal{F}$ with reason ``session limit reached''
    \State \Return \Comment{Maximum 2 hacks per session}
\EndIf
\State \texttt{// DETECT: classify scope}
\State $k \gets \text{ClassifyScope}(\mathcal{F})$ \Comment{Algorithm~\ref{alg:classify}}
\State $\mathcal{R} \gets \text{TargetComponent}(k, \tau)$
\State \texttt{// ISOLATE: create environment}
\If{$k = 1$}
    \State $B \gets \text{InProcessSandbox}()$
\Else
    \State $B \gets \text{GitWorktree}(\mathcal{R}, \texttt{origin/main})$
    \State \textbf{assert} $B.\text{rootDir} \cap \mathcal{F}_{\text{system}} = \emptyset$
\EndIf
\State $\text{fixAttempts} \gets 0$
\Repeat \Comment{MODIFY--TEST loop}
    \State \texttt{// MODIFY: generate change}
    \State $\Delta \gets \text{AgentWrite}(\mathcal{F}, k, B, \text{failureContext})$
    \State \texttt{// TEST: validate}
    \State $\text{result} \gets \Gamma_k(\Delta, \Sigma)$
    \State $\text{fixAttempts} \gets \text{fixAttempts} + 1$
    \If{$\text{result} = \text{fail}$}
        \State $\text{failureContext} \gets \text{ExtractFailureDetails}(\Delta)$
    \EndIf
\Until{$\text{result} = \text{pass}$ \textbf{or} $\text{fixAttempts} \geq 3$}
\If{$\text{result} = \text{fail}$}
    \State \textbf{Cleanup}($B$) \Comment{Remove worktree}
    \State \textbf{Archive} $\mathcal{F}$ with reason ``validation failed after 3 attempts''
    \State \Return
\EndIf
\State \texttt{// GATE: request approval}
\State $\text{decision} \gets \text{GateRequest}(k, \Delta)$
\If{$\text{decision} = \text{rejected}$}
    \State \textbf{Cleanup}($B$)
    \State \textbf{Archive} $\mathcal{F}$ with reason ``gate rejected''
    \State \Return
\EndIf
\State \texttt{// DEPLOY: apply modification}
\State $D_k(\Delta)$ \Comment{Tier-appropriate deployment}
\State \textbf{Cleanup}($B$) \Comment{git worktree remove --force}
\State $H \gets H + 1$
\State \textbf{Record} $(\Delta, k, \tau, \text{timestamp})$ in experience bank
\end{algorithmic}
\end{algorithm}

\begin{algorithm}[H]
\caption{ClassifyScope: Map Friction to Modification Tier}
\label{alg:classify}
\begin{algorithmic}[1]
\Require Friction event $\mathcal{F} = (\tau, e, n)$
\Ensure Tier $k \in \{1, 2, 3, 4\}$
\If{$n < 3$ \textbf{and} $e$ is session-local}
    \State \Return $1$ \Comment{Runtime tool: transient gap}
\ElsIf{$e$ is a missing capability \textbf{and} no existing extension covers $\tau$}
    \State \Return $2$ \Comment{Extension: persistent gap in this agent's harness}
\ElsIf{$e$ is a knowledge or procedure gap \textbf{and} $n > 5$}
    \State \Return $3$ \Comment{Skill: recurring knowledge gap}
\ElsIf{$e$ is a system design gap affecting multiple agents}
    \State \Return $4$ \Comment{Cross-component: ecosystem-level fix needed}
\Else
    \State \Return $1$ \Comment{Default: safest tier}
\EndIf
\end{algorithmic}
\end{algorithm}

\subsection{Session Budget and Time-Boxing}
\label{sec:budget}

To prevent the modification loop from consuming unbounded session resources,
HHP enforces two constraints:

\begin{invariant}[Session Modification Budget]
In any single agent session, the total number of completed modifications
(successful DEPLOY transitions) is bounded by $H_{\max} = 2$.
\end{invariant}

\begin{invariant}[Hack Cycle Time-Box]
Each execution of HARNESS-HACK is bounded by a wall-clock timeout of
$T_{\text{hack}} = 5$ minutes. If the timeout is reached before DEPLOY,
the hack cycle is abandoned and the friction event is archived.
\end{invariant}

These two invariants together ensure that self-modification overhead is bounded:
at most $2 \times 5 = 10$ minutes per session can be spent on harness hacking,
independent of task difficulty or modification complexity.

\begin{theorem}[Session Overhead Bound]
\label{thm:overhead}
Under HHP, the fraction of session time spent on self-modification is at most
$10 / T_{\text{session}}$ where $T_{\text{session}}$ is the total session
duration.
\end{theorem}

\begin{proof}
By Invariants 2 and 3, each session contains at most $H_{\max} = 2$ hack cycles,
each lasting at most $T_{\text{hack}} = 5$ minutes. Total modification time is
bounded by $H_{\max} \cdot T_{\text{hack}} = 10$ minutes. Since
$T_{\text{session}} \geq T_{\text{hack}}$ (a session contains at least one
cycle), the fraction is at most $10 / T_{\text{session}}$. \qed
\end{proof}

% ============================================================================
% 6. INTERACTION WITH CONTINUOUS LEARNING
% ============================================================================
